\section{Safety is Contextual, LLM Judges Are Not: Navigating the Rigid Priors of Evaluators}

Safety labels depend on how intent, assistance, and effect interact in context.
A refusal prefix does not make a response safe when actionable instructions
follow it, while lexical overlap with a harmful request does not make a benign
explanation unsafe. Evaluation therefore requires a joint reading of the user
request, the assistant response, and the criterion being applied.

LLM judges often replace this joint reasoning with stable but incompatible
priors. One evaluator may treat refusal language as decisive, another may focus
on operational detail, and a third may interpret a fictional or mathematical
wrapper as evidence that the content is benign. These priors become visible
when the input is held constant across judges. Our analysis uses two rubrics: a
binary classifier that directly labels response harmfulness and a graded rubric
that combines refusal, convincingness, and specificity. Differences across both
backbones and rubrics reveal whether disagreement originates in the evaluator,
the operationalization of safety, or their interaction.
